\section{Penrose 奇性定理}\label{sec:penrose-singularity-theorem}

设 $(M,g)$ 为整体双曲时空, $S$ 为带有指向未来的单位法向量场 $n$ 的 Cauchy 超曲面, $\Sigma\subset S$ 为在 $S$ 中具有单位法向量场 $\nu$ 的紧二维子流形. 对每个 $p\in\Sigma$, 设 $c_p$ 为初始条件为 $n_p+\nu_p$ 的零模测地线. 对某个 $\varepsilon>0$ 定义光滑映射 $\exp:(-\varepsilon,\varepsilon)\times\Sigma\to M$, 使 $\exp(r,p)=c_p(r)$.

\begin{RiemannDefinition}\label{def:conjugate-points-to-spacelike-surface}
$\exp$ 的临界值称为 $\Sigma$ 的\textbf{共轭点}.
\end{RiemannDefinition}

粗略地说, 共轭点是从 $\Sigma$ 上邻近点正交出发的测地线相交之处 [另见第~\ref{chap:riemannian-manifolds}~章的 Exercise~\ref{ex:minimizing-properties-of-geodesics}(\ref{ex:jacobi-fields-and-cut-locus})].

设 $q=\exp(r_0,p)$ 是不与 $\Sigma$ 共轭的一点. 若 $\varphi$ 是 $\Sigma$ 在 $p$ 附近的局部参数化, 则可利用映射
\[
  (u,r,x^2,x^3)\longmapsto \exp(r,\psi_u(\varphi(x^2,x^3)))
\]
在某个开集 $V\ni q$ 上构造局部坐标系 $(u,r,x^2,x^3)$, 其中 $\psi_u$ 是沿正交于 $S$ 的类时测地线的流, 映射 $\exp:(-\varepsilon,\varepsilon)\times\psi_u(\Sigma)\to M$ 的定义同上.

由于 $\frac{\partial}{\partial r}$ 与零模测地线相切, 有 $g_{rr}=\langle\frac{\partial}{\partial r},\frac{\partial}{\partial r}\rangle=0$. 另一方面, 有
\begin{align*}
  \frac{\partial g_{r\mu}}{\partial r}
  &=\frac{\partial}{\partial r}\left\langle\frac{\partial}{\partial r},
                   \frac{\partial}{\partial x^\mu}\right\rangle
   =\left\langle\frac{\partial}{\partial r},
       \nabla_{\frac{\partial}{\partial r}}\frac{\partial}{\partial x^\mu}\right\rangle\\
  &=\left\langle\frac{\partial}{\partial r},
       \nabla_{\frac{\partial}{\partial x^\mu}}\frac{\partial}{\partial r}\right\rangle
   =\frac12\frac{\partial}{\partial x^\mu}
          \left\langle\frac{\partial}{\partial r},\frac{\partial}{\partial r}\right\rangle=0,
\end{align*}
对 $\mu=0,1,2,3$ 成立. 由于在 $\psi_u(\Sigma)$ 上 $g_{ru}=-1$ 且 $g_{r2}=g_{r3}=0$, 在 $V$ 上也有 $g_{ru}=-1$ 和 $g_{r2}=g_{r3}=0$. 因而度量在此坐标系下写为
\begin{align*}
  g&=\alpha\,\mathrm{d}u\otimes\mathrm{d}u-\mathrm{d}u\otimes\mathrm{d}r
       -\mathrm{d}r\otimes\mathrm{d}u\\
   &\quad+\sum_{i=2}^3\beta_i(\mathrm{d}u\otimes\mathrm{d}x^i+\mathrm{d}x^i\otimes\mathrm{d}u)
       +\sum_{i,j=2}^3\gamma_{ij}\,\mathrm{d}x^i\otimes\mathrm{d}x^j.
\end{align*}
由于
\[
  \det\begin{pmatrix}
    \alpha & -1 & \beta_2 & \beta_3\\
    -1 & 0 & 0 & 0\\
    \beta_2 & 0 & \gamma_{22} & \gamma_{23}\\
    \beta_3 & 0 & \gamma_{32} & \gamma_{33}
  \end{pmatrix}
  =-\det\begin{pmatrix}
    \gamma_{22} & \gamma_{23}\\
    \gamma_{32} & \gamma_{33}
  \end{pmatrix},
\]
可见函数
\[
  \gamma_{ij}:=\left\langle\frac{\partial}{\partial x^i},\frac{\partial}{\partial x^j}\right\rangle
\]
组成一个正定矩阵, 因而 $g$ 在二维曲面 $\exp(r,\psi_u(\Sigma))$ 上诱导 Riemann 度量, 这些曲面因而是类空的. 由于向量场 $\frac{\partial}{\partial x^i}$ 总能沿 $c_p$ 定义, 即便在坐标系失效的点, 即在 $\Sigma$ 的共轭点处, 矩阵 $(\gamma_{ij})$ 也沿 $c_p$ 良定义. 这些正是 $\gamma:=\det(\gamma_{ij})$ 为零的点, 因为只有在这些点处, $\{\frac{\partial}{\partial u},\frac{\partial}{\partial r},\frac{\partial}{\partial x^2},\frac{\partial}{\partial x^3}\}$ 才不再线性无关. (事实上, 向量场 $\frac{\partial}{\partial x^i}$ 是沿 $c_p$ 的 Jacobi 向量场 --- 见第~\ref{chap:riemannian-manifolds}~章的 Exercise~\ref{ex:minimizing-properties-of-geodesics}(\ref{ex:jacobi-fields-and-cut-locus}).)

容易看出
\[
  \Gamma^u_{ur}=\Gamma^u_{rr}=\Gamma^u_{ri}=\Gamma^r_{rr}=\Gamma^i_{rr}=0
  \qquad\text{且}\qquad\Gamma^i_{rj}=\sum_{k=2}^3\gamma^{ik}\beta_{kj},
\]
其中 $(\gamma^{ij})=(\gamma_{ij})^{-1}$, $\beta_{ij}=\frac12\frac{\partial\gamma_{ij}}{\partial r}$ [参见 Exercise~\ref{ex:penrose-singularity-theorem}(\ref{ex:christoffel-symbols-in-null-coordinates})]. 因此
\begin{align*}
  R_{rr}
  &=R_{urr}^{\ \ \ u}+\sum_{i=2}^3R_{irr}^{\ \ \ i}
    =\sum_{i=2}^3\left(-\frac{\partial\Gamma^i_{ir}}{\partial r}
                     -\sum_{j=2}^3\Gamma^j_{ir}\Gamma^i_{rj}\right)\\
  &=-\frac{\partial}{\partial r}\left(\sum_{i,j=2}^3\gamma^{ij}\beta_{ij}\right)
     -\sum_{i,j,k,l=2}^3\gamma^{jk}\gamma^{il}\beta_{ki}\beta_{lj}.
\end{align*}
(参见第~\ref{chap:curvature}~章第~\ref{sec:curvature}~节.) 此式中出现的量
\[
  \theta:=\sum_{i,j=2}^3\gamma^{ij}\beta_{ij}
\]
称为零模测地线的\textbf{膨胀率}, 具有重要的几何意义:
\[
  \theta=\frac{1}{2}\operatorname{tr}\left((\gamma_{ij})^{-1}
          \frac{\partial}{\partial r}(\gamma_{ij})\right)
       =\frac{1}{2}\frac{\partial}{\partial r}\log\gamma
       =\frac{\partial}{\partial r}\log\gamma^{1/2},
\]
其中 $\gamma:=\det(\gamma_{ij})$. 因而, 膨胀率给出了类空二维曲面 $\exp(r,\psi_u(\Sigma))$ 的面积元的变化. 对我们而言更重要的是, 膨胀率的奇点表明 $\gamma$ 取零值, 即出现 $\psi_u(\Sigma)$ 的一个共轭点.

\begin{RiemannDefinition}\label{def:null-energy-condition}
若对任意零模向量场 $V\in\mathfrak{X}(M)$ 均有 $\operatorname{Ric}(V,V)\geq0$, 则称时空 $(M,g)$ 满足\textbf{零模能量条件}.
\end{RiemannDefinition}

容易看出, 强能量条件蕴含这一条件 (但这一条件较弱). 由 Einstein 方程, 这等价于要求约化能量--动量张量 $T$ 对任意零模向量场 $V\in\mathfrak{X}(M)$ 满足 $T(V,V)\geq0$. 对于静止质量密度函数为 $\rho\in C^\infty(M)$, 单位速度向量场为 $U\in\mathfrak{X}(M)$ 的无压流体, 该要求成为
\[
  \rho\langle U,V\rangle^2\geq0\iff\rho\geq0.
\]
对于更复杂的物质模型, 零模能量条件也给出同样合理的限制.

\begin{RiemannProposition}\label{prop:null-focusing-under-null-energy-condition}
设 $(M,g)$ 为满足零模能量条件的整体双曲时空, $S\subset M$ 为 Cauchy 超曲面, $\Sigma\subset S$ 为在 $S$ 中具有单位法向量场 $\nu$ 的紧二维子流形, 且 $p\in\Sigma$ 为满足 $\theta=\theta_0<0$ 的一点. 则零模测地线 $c_p$ 至少包含一个 $\Sigma$ 的共轭点, 距 $\Sigma$ 向未来的仿射参数距离至多为 $-\frac{2}{\theta_0}$ (假定它能延拓至该距离).
\end{RiemannProposition}

\begin{RiemannProof}
由于 $(M,g)$ 满足零模能量条件, 有 $R_{rr}=\operatorname{Ric}(\frac{\partial}{\partial r},\frac{\partial}{\partial r})\geq0$. 因此,
\[
  \frac{\partial\theta}{\partial r}
     +\sum_{i,j,k,l=2}^3\gamma^{jk}\gamma^{il}\beta_{ki}\beta_{lj}\leq0.
\]
选取一个标准正交基 (此时 $\gamma^{ij}=\delta_{ij}$), 并利用 $n\times n$ 方阵满足的不等式
\[
  (\operatorname{tr}A)^2\leq n\operatorname{tr}(A^tA),
\]
容易证明
\[
  \sum_{i,j,k,l=2}^3\gamma^{jk}\gamma^{il}\beta_{ki}\beta_{lj}
    =\sum_{i,j=2}^3\beta_{ji}\beta_{ij}
    =\operatorname{tr}\bigl((\beta_{ij})(\beta_{ij})^t\bigr)\geq\frac12\theta^2.
\]
所以 $\theta$ 必定满足
\[
  \frac{\partial\theta}{\partial r}+\frac12\theta^2\leq0.
\]
对这一不等式积分得到
\[
  \frac{1}{\theta}\geq\frac{1}{\theta_0}+\frac{r}{2},
\]
因而 $\theta$ 必定在 $r$ 不大于 $-\frac{2}{\theta_0}$ 时发散.
\end{RiemannProof}

将紧曲面 $\Sigma$ 的\phantomsection\label{idx:chronological-and-causal-sets-of-compact-surface}\textbf{时序未来}和\textbf{因果未来}定义为
\[
  I^+(\Sigma)=\bigcup_{p\in\Sigma}I^+(p)
  \qquad\text{和}\qquad
  J^+(\Sigma)=\bigcup_{p\in\Sigma}J^+(p)
\]
($\Sigma$ 的\textbf{时序过去}和\textbf{因果过去}也类似定义). 显然, $I^+(\Sigma)$ 是开集的并, 本身为开集, 且有 $J^+(\Sigma)\subset\overline{I^+(\Sigma)}$ 和 $I^+(\Sigma)=\operatorname{int}J^+(\Sigma)$. 另一方面, 容易将 Proposition~\ref{prop:compact-past-intersect-future-domain-of-dependence} (从而 Corollary~\ref{cor:closed-causal-futures-and-compact-diamonds}) 推广到以紧曲面代替点的相应论断 [参见 Exercise~\ref{ex:penrose-singularity-theorem}(\ref{ex:causal-past-of-compact-surface})]. 特别地, $J^+(\Sigma)$ 是闭集. 因此
\[
  \partial J^+(\Sigma)=\partial I^+(\Sigma)=J^+(\Sigma)\setminus I^+(\Sigma),
\]
所以直接推广 Corollary~\ref{cor:global-boundary-causal-curves-are-null-geodesics} 可知, 此边界上的每点都可由一条朝向未来的零模测地线从 $\Sigma$ 中一点到达. 此外, 这条测地线必须与 $\Sigma$ 正交. 事实上, 在 $\Sigma$ 上有
\[
  \frac{\partial}{\partial u}=n
  \qquad\text{且}\qquad\frac{\partial}{\partial r}=n+\nu,
\]
因而度量具有形式
\[
  g=-\mathrm{d}u\otimes\mathrm{d}u-\mathrm{d}u\otimes\mathrm{d}r
    -\mathrm{d}r\otimes\mathrm{d}u
    +\sum_{i,j=2}^3\gamma_{ij}\,\mathrm{d}x^i\otimes\mathrm{d}x^j.
\]
若 $c:I\subset\mathbb{R}\to M$ 是满足 $c(0)\in\Sigma$ 的朝向未来的零模测地线, 则它的初始切向量
\[
  \dot c(0)=\dot u\frac{\partial}{\partial u}
         +\dot r\frac{\partial}{\partial r}
         +\sum_{i=2}^3\dot x^i\frac{\partial}{\partial x^i}
       =(\dot u+\dot r)n+\dot r\nu
         +\sum_{i=2}^3\dot x^i\frac{\partial}{\partial x^i}
\]
必须满足
\[
  \dot u(\dot u+2\dot r)=\sum_{i,j=2}^3\gamma_{ij}\dot x^i\dot x^j.
\]
由于 $c$ 朝向未来, 必有 $\dot u+\dot r>0$. 另一方面, 在 $S$ 中可选取 $\nu$ 或 $-\nu$ 为 $\Sigma$ 的单位法向量, 所以可以假设 $\dot r\geq0$. 若 $c$ 不与 $\Sigma$ 正交, 则有
\[
  \sum_{i,j=2}^3\gamma_{ij}\dot x^i\dot x^j>0
    \implies\dot u(\dot u+2\dot r)>0\implies\dot u>0.
\]
现在 $u>0$ 且 $r\geq0$ 的区域显然是 $I^+(\Sigma)$ 的子集, 因为其中各点都可由一段类时测地线与一段零模测地线组成的分段光滑曲线从 $\Sigma$ 到达. 因此可见, 若 $c$ 不与 $\Sigma$ 正交, 则对所有 $t>0$ 都有 $c(t)\in I^+(\Sigma)$.

即使是正交于 $\Sigma$ 的朝向未来的零模测地线, 最终也可能进入 $I^+(\Sigma)$. 下列结果给出发生这种情况的一个充分条件.

\begin{RiemannProposition}\label{prop:null-geodesic-enters-chronological-future-past-conjugate-point}
设 $(M,g)$ 为整体双曲时空, $S$ 为带有指向未来的单位法向量场 $n$ 的 Cauchy 超曲面, $\Sigma\subset S$ 是在 $S$ 中带有单位法向量场 $\nu$ 的紧二维子流形, $p\in\Sigma$, $c_p$ 是过 $p$ 且初始条件为 $n_p+\nu_p$ 的零模测地线, 且对某个 $r>0$ 有 $q=c_p(r)$. 若 $c_p$ 在 $p$ 与 $q$ 之间有一个共轭点, 则 $q\in I^+(\Sigma)$.
\end{RiemannProposition}

\begin{RiemannProof}
只给出证明的概要. 设 $s$ 为沿 $c_p$ 在 $p$ 与 $q$ 之间的第一个共轭点. 由于 $q$ 与 $p$ 共轭, 存在另一条从 $\Sigma$ 出发的零模测地线 $\gamma$, 在 $s$ 处 (近似) 与 $c_p$ 相交. 先沿 $\gamma$ 从 $\Sigma$ 到 $s$, 再沿 $c_p$ 从 $s$ 到 $q$, 所得分段光滑的零模曲线是因果曲线, 但不是零模测地线. 容易将此曲线光滑化, 同时保持它为因果曲线但不成为测地线. 因而由 Corollary~\ref{cor:global-boundary-causal-curves-are-null-geodesics} 的推广可得 $q\in I^+(\Sigma)$ (见 Fig.~\ref{fig:null-geodesic-past-conjugate-point}).
\end{RiemannProof}
\RGfigure[.62\textwidth]{null-geodesic-past-conjugate-point}{Proposition~\ref{prop:null-geodesic-enters-chronological-future-past-conjugate-point} 的证明}

\begin{RiemannDefinition}\label{def:trapped-surface}
设 $(M,g)$ 为整体双曲时空, $S$ 为带有指向未来的单位法向量场 $n$ 的 Cauchy 超曲面. 设 $\Sigma\subset S$ 为在 $S$ 中带有单位法向量场 $\nu$ 的紧二维子流形. 若初始条件为 $n+\nu$ 和 $n-\nu$ 的零模测地线的膨胀率 $\theta^+$ 和 $\theta^-$ 在 $\Sigma$ 上处处均为负, 则称 $\Sigma$ 为\textbf{陷俘曲面}.
\end{RiemannDefinition}

现在已经具备证明 Penrose 奇性定理的所有条件.

\begin{RiemannTheorem}[Penrose]\label{thm:penrose-null-singularity}
设 $(M,g)$ 为带有非紧 Cauchy 超曲面 $S$ 的连通整体双曲时空, 且满足零模能量条件. 若 $S$ 含有一个陷俘曲面 $\Sigma$, 则 $(M,g)$ 是奇性时空.
\end{RiemannTheorem}

\begin{RiemannProof}
设 $t:M\to\mathbb{R}$ 是使 $S=t^{-1}(0)$ 的全局时间函数. $\operatorname{grad}t$ 的积分曲线为类时曲线, 与 $S$ 恰好相交一次, 且与 $\partial I^+(\Sigma)$ 至多相交一次. 这就定义了一个连续单射 $\pi:\partial I^+(\Sigma)\to S$, 其像是开集. 事实上, 若 $q=\pi(p)$, 则 $q$ 的某个邻域内的所有点都是 $\partial I^+(\Sigma)$ 中某些点的像. 否则存在 $S$ 中的点列 $q_n\to q$, 使 $\operatorname{grad}t$ 经过 $q_n$ 的积分曲线不与 $\partial I^+(\Sigma)$ 相交. 沿 $\operatorname{grad}t$ 的积分曲线取 $p$ 未来的某点 $r$, 令 $r_n$ 为这些曲线与 Cauchy 超曲面 $t^{-1}(t(r))$ 的交点. 则 $r_n\to r$, 因而当 $n$ 足够大时 $r_n\in I^+(\Sigma)$ (因为 $I^+(\Sigma)$ 是开集), 从而得到矛盾.

由于 $\Sigma$ 是陷俘的 (且紧), 存在 $\theta_0<0$, 使与 $\Sigma$ 正交的零模测地线的膨胀率 $\theta^+$ 和 $\theta^-$ 都满足 $\theta^+,\theta^-\leq\theta_0$. 将证明存在一条与 $\Sigma$ 正交且朝向未来的零模测地线, 它在 $\Sigma$ 的未来不能延拓至大于 $r_0=-\frac{2}{\theta_0}$ 的仿射参数. 假设不然. 根据 Proposition~\ref{prop:null-focusing-under-null-energy-condition}, 任意与 $\Sigma$ 正交的零模测地线都将在 $\Sigma$ 的未来至多仿射参数距离 $r_0$ 处出现共轭点; 越过该点后, 由 Proposition~\ref{prop:null-geodesic-enters-chronological-future-past-conjugate-point}, 它将处于 $I^+(\Sigma)$ 中. 因此, $\partial I^+(\Sigma)$ 是紧集
\[
  \exp^+([0,r_0]\times\Sigma)\cup\exp^-([0,r_0]\times\Sigma)
\]
的一个 (闭) 子集 ($\exp^+$ 和 $\exp^-$ 分别表示用单位法向量 $\nu$ 和 $-\nu$ 构造的指数映射), 因而是紧的. 所以 $\pi$ 的像也紧, 因而既闭又开. 由于 $M$, 从而 $S$, 是连通的, $\pi$ 的像将为 $S$, 此时 $S$ 将与 $\partial I^+(\Sigma)$ 同胚. 但依假设 $S$ 非紧, 这就得出矛盾.
\end{RiemannProof}

\begin{RiemannRemark}\label{rem:anti-trapped-surface}
显然, 若把存在陷俘曲面的条件换为存在\textbf{反陷俘曲面}的条件, 则 $(M,g)$ 仍是奇性时空. 反陷俘曲面是紧曲面 $\Sigma\subset S$, 使与 $\Sigma$ 正交的零模测地线的两个膨胀率均为正. 在此情形, 存在一条与 $\Sigma$ 正交且\textbf{朝向过去}的零模测地线, 它在 $\Sigma$ 的\textbf{过去}不能延拓至大于 $r_0=\frac{2}{\theta_0}$ 的仿射参数时间.
\end{RiemannRemark}

\begin{RiemannExample}\label{ex:penrose-singularities-in-schwarzschild-and-flrw}
\begin{enumerate}[label=(\arabic*),ref=\arabic*,leftmargin=*]
\item Schwarzschild 解的区域 $\{r<2m\}$ 整体双曲 [参见 Exercise~\ref{ex:causality}(\ref{ex:globally-hyperbolic-spacetime-examples})], 且满足零模能量条件 (因为 $\operatorname{Ric}=0$). 由于 $r$ (或 $-r$) 显然是时间函数 (取决于时间定向的选择), 它必沿任意朝向未来的零模测地线增大 (或减小), 所以任意 $(t,r)$ 为常数的球面 $\Sigma$ 都是反陷俘 (或陷俘) 的. 由于任意 Cauchy 超曲面都与 $\mathbb{R}\times S^2$ 微分同胚, 因而非紧, 由 Theorem~\ref{thm:penrose-null-singularity} 可知 Schwarzschild 解在 $\Sigma$ 的过去 (或未来) 有奇性. 此外, Theorem~\ref{thm:hawking-timelike-singularity} 蕴含这一奇性是普遍的: 满足零模能量条件的 Schwarzschild 解的任意足够小的扰动, 也将有奇性. 粗略地说, 一旦坍缩进行得足够久, 就无法阻止奇点形成.
\item FLRW 模型整体双曲 [参见 Exercise~\ref{ex:causality}(\ref{ex:globally-hyperbolic-spacetime-examples})], 且满足零模能量条件 (因为 $\rho>0$). 此外, 径向零模测地线满足
\[
  \frac{\mathrm{d}r}{\mathrm{d}t}=\pm\frac{1}{a}\sqrt{1-kr^2}.
\]
因此, 若从 $(t,r)$ 为常数的一个球面 $\Sigma$ 出发, 沿 $r$ 增大或减小的方向追踪与其正交的零模测地线, 则得到半径 $ar$ 满足
\[
  \frac{\mathrm{d}}{\mathrm{d}t}(ar)=\dot a r+a\dot r=\dot a r\pm\sqrt{1-kr^2}
\]
的球面. 假定模型正在膨胀, 大爆炸发生于 $t=0$, 且空间非紧 (特别地, $k\ne1$). 则对足够小的 $t>0$, 球面 $\Sigma$ 是反陷俘的, 因此 Theorem~\ref{thm:penrose-null-singularity} 保证该模型在 $\Sigma$ 的过去有奇性 (即存在大爆炸). 此外, Theorem~\ref{thm:penrose-null-singularity} 蕴含这一奇性是普遍的: 一个正在膨胀, 空间非紧且满足零模能量条件的 FLRW 模型的任意足够小的扰动, 也将有奇性. 粗略地说, 任何膨胀宇宙都必从大爆炸开始.
\end{enumerate}
\end{RiemannExample}

\begin{RiemannExercise}\label{ex:penrose-singularity-theorem}
\begin{enumerate}[label=(\arabic*),ref=\arabic*,leftmargin=*]
\item\label{ex:christoffel-symbols-in-null-coordinates}
证明度量
\begin{align*}
  g&=\alpha\,\mathrm{d}u\otimes\mathrm{d}u-\mathrm{d}u\otimes\mathrm{d}r
       -\mathrm{d}r\otimes\mathrm{d}u
       +\sum_{i=2}^3\beta_i(\mathrm{d}u\otimes\mathrm{d}x^i+\mathrm{d}x^i\otimes\mathrm{d}u)\\
   &\quad+\sum_{i,j=2}^3\gamma_{ij}\,\mathrm{d}x^i\otimes\mathrm{d}x^j
\end{align*}
的 Christoffel 符号满足
\[
  \Gamma^u_{ur}=\Gamma^u_{rr}=\Gamma^u_{ri}=\Gamma^r_{rr}=\Gamma^i_{ur}=\Gamma^i_{rr}=0
  \qquad\text{且}\qquad
  \Gamma^i_{rj}=\sum_{k=2}^3\gamma^{ik}\beta_{kj},
\]
其中 $(\gamma^{ij})=(\gamma_{ij})^{-1}$, $\beta_{ij}=\frac12\frac{\partial\gamma_{ij}}{\partial r}$.
\item\label{ex:causal-past-of-compact-surface}
设 $(M,g)$ 为整体双曲时空, Cauchy 超曲面 $S_0$ 和 $S_1$ 满足 $S_1\subset D^+(S_0)$, $\Sigma\subset S_1$ 为紧曲面. 证明:
\begin{enumerate}[label=(\alph*),leftmargin=*]
\item $D^+(S_0)\cap J^-(\Sigma)$ 是紧集;
\item $J^-(\Sigma)$ 是闭集.
\end{enumerate}
\item\label{ex:limits-of-penrose-singularity-theorem}
解释为何 Penrose 奇性定理不适用于下列各个时空:
\begin{enumerate}[label=(\alph*),leftmargin=*]
\item Minkowski 时空;
\item Einstein 宇宙 [参见 Exercise~\ref{ex:cosmology}(\ref{ex:einstein-and-de-sitter-universe-embeddings})];
\item de Sitter 宇宙 [参见 Exercise~\ref{ex:cosmology}(\ref{ex:einstein-and-de-sitter-universe-embeddings})];
\item 二维反 de Sitter 宇宙 [参见 Exercise~\ref{ex:hawking-singularity-theorem}(\ref{ex:anti-de-sitter-geodesic-completeness})].
\end{enumerate}
\end{enumerate}
\end{RiemannExercise}
